% 2026-11-16 站主問：「賽局理論，本文如何理解運用，而本文的假設又如何有其欠缺不足？」「他的論述都有很多假設在裡面」「作者的盲點」。
% 引句由兩路 agent 自 readings/W10/text/ 文字層抽出，關鍵十餘句主對話回文字層核過頁碼（PDF 頁 + 722）。
\chapter{賽局理論之運用、假設與盲點}
\label{ch:20_game}

\begin{qaquestion}
Law\ix{persons}{law@Law, David S.} 以賽局理論說明無兵無財的法院何以獲得服從，並據以主張司法審查支撐人民主權。
\begin{enumerate}
\tightlist
\item 本文用了哪幾個模型？
\item 各模型以何種假設為前提？
\item 這些假設在何處不成立？
\item 模型未納入、而會改變結論的因素有哪些？
\end{enumerate}
\end{qaquestion}

\section{本文運用的模型}
\label{sec:20_game-1}

\begin{enumerate}
\tightlist
\item
  委託人與代理人\ix{terms}{01@\zhlidxcat{賽局理論}!委託人與代理人}。Law 以 Hobbestown 的寓言起頭：居民把槍交給一名槍手，人民整體為委託人，槍手為代理人\pg{737 頁}。代理人可能把受託的武力轉向委託人，Kiewiet\ix{persons}{kiewiet@Kiewiet, D. Roderick} 與 McCubbins\ix{persons}{mccubbins@McCubbins, Mathew D.} 稱之為 Madison's dilemma\ix{terms}{01@\zhlidxcat{賽局理論}!Madison's dilemma}\pg{737 頁}。人民控制槍手有三重障礙：授權條款與違約對策須先議定，而議定於無知之幕後，契約必不完整\pg{739 頁}；人民不知槍手是否越權\pg{740 頁}；推翻槍手須集體行動，個別起事者必死\pg{741 頁}。\footcite[737--741]{law2009}
\item
  協調賽局\ix{terms}{01@\zhlidxcat{賽局理論}!協調賽局}與共同知識\ix{terms}{01@\zhlidxcat{賽局理論}!共同知識}。推翻槍手的條件是人人知道槍手越權、人人知道人人知道，依此遞推\pg{742 頁}。人人不動也是均衡，Law 在註 57 引 Sutter\ix{persons}{sutter@Sutter, Daniel}：「Universal nonparticipation may always be an equilibrium」〔全體不參與始終可能是一個均衡〕\pg{742 頁}。因此人民需要 Schelling\ix{persons}{schelling@Schelling, Thomas C.} 所說的協調信號。\footcite[742]{law2009}
\item
  焦點（focal point\ix{terms}{01@\zhlidxcat{賽局理論}!焦點}）。法院宣告政府越權，就是這個信號；判決使人相信反抗既有必要也會有人跟進，起事於是成為自我實現的預言\pg{744 頁}。Law 以紅綠燈為喻：路口各車目標相同，只需一個人人看得見、人人知道別人也看得見的信號\pg{770 頁}；焦點須顯著且唯一，終審憲法法院兩者兼具\pg{771、774--775 頁}。\footcite[744, 770--775]{law2009}
\item
  均衡\ix{terms}{01@\zhlidxcat{賽局理論}!均衡}與聲譽。邊界爭端兩國都寧可讓步也不願開戰，法院判決選定其中一個均衡，敗訴國抗拒即招致戰爭\pg{759--760 頁}。法院的權力由此取決於判決是否被遵守，與判決受不受歡迎無關：「one need not be popular in order to be powerful」〔有權力不必受歡迎〕\pg{780 頁}，Bush v. Gore\ix{cases}{02@美國!01@聯邦最高法院!2000Bush v. Gore（2000）@Bush v. Gore（2000）} 因此被計為法院權力的增加\pg{782 頁}。反過來，法院不能讓人看見不服從而無後果的偏離均衡行為\pg{781 頁}。\footcite[759--760, 780--782]{law2009}
\item
  第五部的訊號賽局\ix{terms}{01@\zhlidxcat{賽局理論}!訊號賽局}。參與者為政府、法院、人民；人民先決定是否讓法院獨立，政府再決定是否違憲，法院判政府勝敗\pg{787 頁}。Law 假設政府與法院對事實與法律有完全資訊，人民兩者皆無，也無從判斷法院的訊號是否正確\pg{787--788 頁}。結論是司法獨立\ix{terms}{02@\zhlidxcat{司法權之基礎}!司法獨立}提高的是正面訊號（判政府違憲）的比率，正確訊號的比率未必提高：誤判違憲增加，誤判合憲減少\pg{788--789 頁}。人民厭惡暴政的風險，因而寧可忍受過度積極的法院，理由是法院自己沒有壓迫人民的手段\pg{789--790 頁}。\footcite[787--790]{law2009}
\end{enumerate}

Law 拿這套模型回應抗多數難題\ix{terms}{02@\zhlidxcat{司法權之基礎}!抗多數難題}：司法審查使人民得以監督與協調，屬於人民主權\ix{terms}{02@\zhlidxcat{司法權之基礎}!人民主權}的工具\pg{786 頁}。他對正當性理論的批評也從這裡出發：那套理論迴避了官員服從法院究竟出於規範上的義務還是自利\pg{793 頁}。\footcite[786, 793]{law2009}

\section{本文的假設}
\label{sec:20_game-2}

\begin{longtable}{@{}L{0.50\linewidth}L{0.22\linewidth}L{0.22\linewidth}@{}}
\caption{Law 一文各模型所依據的假設。}\\
\toprule
假設 & 頁次 & 所用的模型 \\
\midrule
\endfirsthead
\toprule
假設 & 頁次 & 所用的模型 \\
\midrule
\endhead
人民為單一的集體委託人，偏好推翻濫權的槍手 & 731、737、741 & 委託代理、協調 \\
行為者理性，依對他人行為的預期行動 & 757、758 & 協調、焦點 \\
訴訟當事人與公民的共同利益大於衝突 & 758--759、763 & 均衡選擇 \\
憲法為授權契約，政府行為違反契約與否有確定答案 & 738、787 & 監督、訊號賽局 \\
政府與法院資訊完全，人民無資訊 & 787 & 訊號賽局 \\
獨立的法院具專業且無欺騙誘因 & 786 及註 204 & 訊號可信 \\
法院自己無壓迫人民的手段 & 789--790 & 人民容忍誤判違憲 \\
終極制裁為全民起事 & 741--742、766--769 & 可信威脅 \\
\bottomrule
\end{longtable}

\section{假設的欠缺}
\label{sec:20_game-3}

\begin{enumerate}
\tightlist
\item
  第五部的賽局假設政府行為合憲與否有確定答案，Law 自承「the setup of the game deliberately ignores the fact that government actions often fall in a grey area of uncertain constitutionality, and that courts cannot merely apply the constitution but must instead make law」〔此一賽局的設定刻意忽略政府行為常落於合憲性不明的灰色地帶，而法院不能只是適用憲法，須造法始能裁判此類爭議〕\pg{790 頁}。他在 777 頁又說，難以判斷的案件正是終審憲法法院的日常（close cases are the stock and trade）。抗多數難題的爭點集中在這類案件：墮胎、選舉區劃、平權措施的判決，所涉為憲法條款的內容由誰決定，政府是否越過一條既定界線並無確定答案。模型排除了灰色地帶，也就排除了抗多數難題發生的那一類案件。739 頁承認人民的憲法必有漏洞，填補漏洞者正是法院；法院在此兼為代理人，Madison's dilemma 於法院再次出現。\footcite[739, 777, 790]{law2009}
\item
  Law 在 739 頁引 Weingast\ix{persons}{weingast@Weingast, Barry R.}（註 47）寫道，槍手只需部分人民支持，便可壓迫其餘人民並與支持者分贓；單一集體委託人的假設，在同一頁已不成立。人民一旦分裂，推翻槍手對一部分人有利、對另一部分人有害，賽局的結構隨之改變：各方利益分歧，紅綠燈喻所預設的「They all share the same goal」〔他們目標相同〕\pg{770 頁}不再成立。法院判政府違憲，在支持政府的那一部分人看來是黨派的判斷，焦點只對一個陣營有效。\footcite[739, 770]{law2009}
\item
  George v. Albert 一例中，Law 說兩造寧可由法院決定總統，也不願由國會決定或沒有總統，故共同利益「happen to be relatively strong」〔恰好相當強〕\pg{763 頁}。這句話描述的是 Bush v. Gore 事後和平收場的結果，事前無從檢驗；高爭議案件中共同利益是否足夠強，正是判決能否發揮協調作用的分界，本文未提出判準。\footcite[763]{law2009}
\item
  依 788--789 頁，獨立的法院多發出正面訊號，誤判違憲的成本由人民承擔；Law 以法院無兵可用為由認為這些成本「more easily endured」〔較易承受〕\pg{789--790 頁}，第五部的結論即以法院無欺騙誘因、無壓迫手段為前提。法院宣告法律違憲即排除民選機關的政策選擇，其分配效果與立法相同，這一項成本未進入人民的權衡。法院自身的偏好在模型中也付之闕如：法官的意識形態、任命機制、對自身權力的利益，都可能使訊號偏向。Law 承認賽局「does not allow for the possibility that the people have access to other sources of information that would enable them to monitor the behavior of both the government and the court」〔未考慮人民另有資訊來源，得以同時監督政府與法院〕\pg{790 頁}，誰監督法院這個問題，模型未處理。\footcite[788--790]{law2009}
\item
  Law 自承革命是粗糙而昂貴的手段\pg{769 頁}，也承認多數違憲判決之後，人民最合理的反應是什麼都不做\pg{777 頁}；以全民起事為終極制裁，可信威脅只在極端情形成立。日常的服從便須以重複賽局中的聲譽解釋\pg{781 頁}，但本文未把重複賽局模型化，判決與選舉、輿論之間怎麼傳導也停在文字層次。777 頁還指出，人民要靠媒體、學界、政界解讀判決，眾聲喧嘩使協調問題重新出現（a Babel of voices）；Law 對此的回應是法院「if it so chooses」〔若有意為之〕可以寫得一致、簡短、清楚\pg{778 頁}，這是對法院的建議，模型本身沒有推出這個結論。\footcite[769, 777--778, 781]{law2009}
\item
  依 780--782 頁，服從越多則權力越大，權力越大則服從越多；這套循環只能說明一個已經存在的服從均衡如何維持。第一個均衡從哪裡來，Law 在 799 頁說沒有明顯或萬無一失的方法；他對 Marbury\ix{cases}{02@美國!01@聯邦最高法院!1803Marbury v. Madison（1803）@Marbury v. Madison（1803）} 的解讀是法院在不會遭到抗拒時才主張權力，「With the passage of time, the illusion of power has become the reality」〔時間一久，權力的幻象成了現實〕\pg{785 頁}。全民不動也是均衡\pg{742 頁}，那麼哪些法院成功建立服從均衡、哪些失敗，本文只舉成功的例子。\footcite[742, 780--785, 799]{law2009}
\end{enumerate}

\section{作者的盲點}
\label{sec:20_game-4}

\begin{enumerate}
\tightlist
\item
  遭俘虜的法院替政府協調。焦點的效果與判決的方向無關：George v. Albert 一例中，法院不論判決有無說服力，都可以讓全體公民協調在它所選的結果上\pg{764 頁}。法院被政府以任命、增額、改組控制之後，仍保有顯著性與唯一性，其合憲判決便替政府協調出人民的默從。796 頁給暴君的教訓只列兩種手段：正面攻擊司法獨立，以及設立管轄重疊的多個法院使訊號互相矛盾。保留單一法院而換掉其中的法官，不在其列。賽局中司法獨立是人民所選的變數\pg{787 頁}，政府改造法院的行為不在賽局的選項之中。\footcite[764, 787, 796]{law2009}
\item
  人民偏好分歧與憲法上的合理異見\ix{terms}{02@\zhlidxcat{司法權之基礎}!合理異見}。Law 的模型把政府越權當成可以觀察的事實，W06 Waldron\ix{persons}{waldron@Waldron, Jeremy} 的出發點則是公民對權利內容存有合理的歧見。\footcite{waldron2006}依 Waldron 的設定，法院判決表達的是多種合理見解中的一種，其所生的協調是在一個有爭議的答案上協調。777 頁「How Courts Are Not Like Traffic Lights」一節列出判決須解釋、無反應規則、須靠中介等差異，未列出人民對判決結論本身意見相左。\footcite[776--778]{law2009}
\item
  協調功能與規範上的正當性。錯誤的判決只要人人預期他人遵守，同樣構成穩定的均衡；協調理論說明人民何以服從，未處理人民應否服從。Law 在 793 頁批評正當性理論迴避服從的動機問題，他自己的理論則把正當性問題換成了服從問題。W09 Alexander\ix{persons}{alexander@Alexander, Larry} 與 Schauer\ix{persons}{schauer@Schauer, Frederick} 以定紛止爭說明服從最高法院解釋的理由\pg{457--463 頁}，\footcite[457--463]{alexander2000}Law 的協調功能與定紛止爭的利益相近，兩者都以一個人人遵從的答案為價值，不問該答案是否正確。
\item
  例證的單位不一。國際法院的邊界案、歐洲法院的 Cassis de Dijon\ix{cases}{03@歐洲共同體!01@歐洲法院!1979Cassis de Dijon（1979）@Cassis de Dijon（1979）}、美國最高法院的 Bush v. Gore、巴基斯坦的 Musharraf 事件，參與者分別是國家、成員國、選民與政黨、律師與軍方\pg{759--762、791 頁}，法院與執行者的關係各不相同。一套模型涵蓋這些情境，賴以成立的是判決作為公開信號這一共同點，各情境中違抗判決的成本由誰承擔、如何計算，本文未逐一說明。\footcite[759--762, 791]{law2009}
\end{enumerate}

\section{課前可提的問題}
\label{sec:20_game-5}

\begin{enumerate}
\tightlist
\item 790 頁的賽局排除了灰色地帶，而抗多數難題集中在灰色地帶。在灰色地帶的案件，法院的協調功能還能不能支持司法審查的民主性？
\item 法院遭到增額或改組之後，焦點效果是否轉為替政府協調默從？本文第五部對暴君的教訓何以未列此項？
\item 協調的價值與 Alexander 與 Schauer 的定紛止爭有何不同？兩者都無法區分正確與錯誤的判決，Law 對此有無回應？
\end{enumerate}
